\documentclass[11pt]{article}
\usepackage[margin=1in]{geometry}
\usepackage{amsmath,amssymb,amsthm}
\usepackage{enumitem}
\usepackage{graphicx}
\usepackage[hidelinks]{hyperref}

\title{Weak Convergence of TEPROG Without Global Strict Convexity}
\author{Candidate full solution for the FA/Banach discovery run}
\date{22 July 2026}

\newtheorem*{theorem}{Theorem}
\newtheorem*{proposition}{Proposition}
\newtheorem*{remark}{Remark}
\newcommand{\MIN}{\operatorname{MIN}}
\newcommand{\domm}{\operatorname{dom}}
\newcommand{\R}{\mathbb R}

\begin{document}
\maketitle

\section*{Source and target}

The source is Reem, Reich, and De Pierro,
\emph{A Telescoping Bregmanian Proximal Gradient Method Without the Global
Lipschitz Continuity Assumption}, arXiv:1804.10273v4 \cite{RRD}.
Remark 6.4, spanning pages 11--12, asks whether the weak convergence result
continues to hold when the Bregman generator (b) is no longer strictly
convex. The source PDF is included as \texttt{source\_paper.pdf}. The two
following full-width crops show the complete remark across the page break.

\begin{center}
\includegraphics[width=\textwidth]{figures/open_problem_crop_page11.png}

\vspace{0.4em}
\includegraphics[width=\textwidth]{figures/open_problem_crop_page12.png}
\end{center}

The core question is whether ``the weak convergence result holds when \(b\)
is no longer strictly convex.'' The natural literal reading retains the other
hypotheses of Section 5 and Theorem 6.1, in particular Assumption 5.2: the
nested telescopic sets \(S_k\) cover \(C\), and \(b\) is strongly convex on
each \(S_k\), with parameter \(\mu_k>0\).

\section*{Status and novelty check}

This packet gives a candidate full affirmative answer under the literal
retained hypotheses. It also gives a finite-dimensional counterexample to the
broader interpretation in which local strong convexity on the telescopic sets
is removed together with global strict convexity.

A bounded search on 22 July 2026 used the exact paper title, arXiv id, the
open-question phrase, \texttt{TEPROG weak convergence strict convexity}, and
close variants. The 13 works in the OpenAlex citation list for DOI
\texttt{10.1007/s10957-019-01509-8} were screened by title and available
abstract or metadata. No later paper claiming this exact answer was found.
This is negative search evidence, not a comprehensive novelty guarantee, so
novelty confidence is moderate.

\section*{Framework}

Let \(X\) be a reflexive Banach space. Let
\(b:X\to(-\infty,\infty]\) be proper, lower semicontinuous, and convex, and
Gateaux differentiable on \(U=\operatorname{int}(\domm b)\neq\varnothing\).
We do \emph{not} assume that \(b\) is strictly convex on all of \(\domm b\).
For \(x\in\domm b\) and \(y\in U\), set
\[
 B(x,y)=b(x)-b(y)-\langle b'(y),x-y\rangle.
\]
Convexity implies \(B(x,y)\geq0\).

The remaining notation and assumptions are those of Section 5 of the source.
In particular, \(C\) is closed and convex, \(\MIN(F)\neq\varnothing\) is
contained in \(U\), and
\[
 S_1\subseteq S_2\subseteq\cdots\subseteq C,
 \qquad \bigcup_{k\geq1}S_k=C,
\]
where \(b\) is strongly convex on \(S_k\) with parameter \(\mu_k>0\).
The step-size condition (20), the limiting difference property, and bounded
first-type Bregman level sets are assumed exactly as in the weak-convergence
part of Theorem 6.1.

\section*{Proof intuition}

The global strict-convexity axiom is used in the source proof only to show
that two weak cluster points cannot differ. But cluster points are not
arbitrary points of \(\domm b\): the preceding part of the proof puts them in
\(\MIN(F)\subset C\). Since the nested \(S_k\)'s cover \(C\), any two cluster
points lie together in one \(S_N\). Strong convexity on that single set gives
all the separation needed.

More explicitly, the source proof already gives a finite limit
\(\ell_q=\lim_k B(q,x_k)\) for every weak cluster point \(q\). For two cluster
points \(q_1,q_2\), the limiting difference property computes the two directed
divergences as opposite real numbers. Both are nonnegative, so both are zero.
Local strong convexity on a common \(S_N\) then turns zero divergence into
\(q_1=q_2\). This is the whole missing step.

\section*{A localized cluster-point lemma}

\begin{proposition}
Let \((x_k)\) be a bounded sequence in \(U\). Suppose every weak cluster point
belongs to a set \(M\subset U\), the limit
\(
\ell_q=\lim_{k\to\infty}B(q,x_k)
\)
exists and is finite for every weak cluster point \(q\), and \(B\) has the
limiting difference property. If for every two points \(q_1,q_2\in M\) there
is a convex set \(S\subset\domm b\) containing them on which \(b\) is strictly
convex, then \((x_k)\) converges weakly.
\end{proposition}

\begin{proof}
Let \(q_1\) and \(q_2\) be weak cluster points, obtained along subsequences
\(N_1\) and \(N_2\), respectively. The limiting difference property gives
\begin{align*}
B(q_2,q_1)
 &=\lim_{\substack{k\to\infty\\k\in N_1}}
   \bigl(B(q_2,x_k)-B(q_1,x_k)\bigr)
 =\ell_{q_2}-\ell_{q_1},\\
B(q_1,q_2)
 &=\lim_{\substack{k\to\infty\\k\in N_2}}
   \bigl(B(q_1,x_k)-B(q_2,x_k)\bigr)
 =\ell_{q_1}-\ell_{q_2}.
\end{align*}
Thus \(B(q_2,q_1)=-B(q_1,q_2)\). Both terms are nonnegative by convexity,
so both vanish. Strict convexity on a convex set containing \(q_1,q_2\),
together with differentiability at the second argument, implies that zero
Bregman divergence is possible only when \(q_1=q_2\). Hence all weak cluster
points coincide.

It remains to upgrade this uniqueness to convergence. If the sequence did not
converge weakly to its unique cluster point \(q\), some weak neighborhood of
\(q\) would omit an entire subsequence. That bounded subsequence has a weakly
convergent further subsequence by reflexivity, whose limit must again be \(q\),
a contradiction. Therefore \(x_k\rightharpoonup q\).
\end{proof}

\section*{Main theorem: affirmative answer under the stated hypotheses}

\begin{theorem}
Retain the framework and assumptions of the weak-convergence part of Theorem
6.1 of \cite{RRD}, including Assumptions 5.1--5.4 and, in particular, strong
convexity of \(b\) on every telescopic set \(S_k\). Replace only the global
strict-convexity requirement on \(\domm b\) by convexity. Then every TEPROG
sequence generated by Algorithm 5.5 or 5.6 converges weakly to a point of
\(\MIN(F)\). The same conclusion holds under the ``in particular'' sufficient
conditions in Theorem 6.1.
\end{theorem}

\begin{proof}
First note that the proximal subproblem on \(S_k\) is still uniquely solvable:
strong convexity of \(b\) on \(S_k\) implies strict convexity there, so the
objective defining \(p_{L_k,\mu_k,S_k}\) is strictly convex. Thus deleting
global strict convexity creates no selection ambiguity under Assumption 5.2.

All estimates in the proof of Theorem 6.1 through the construction of weak
cluster points remain unchanged. They use convexity of \(b\), strong convexity
on the active sets \(S_k\), and the other retained assumptions, but not strict
convexity outside \(C\). Concretely, those estimates establish:
\begin{enumerate}[label=(\roman*),leftmargin=*]
\item \((x_k)\) is bounded;
\item every weak cluster point \(q\) belongs to \(\MIN(F)\subset C\cap U\);
\item for every such \(q\), the finite limit
      \(\ell_q=\lim_k B(q,x_k)\) exists.
\end{enumerate}

Apply the preceding proposition with \(M=\MIN(F)\). It remains only to verify
its local strictness condition. Given \(q_1,q_2\in M\subset C\), choose
indices \(i,j\) with \(q_1\in S_i\) and \(q_2\in S_j\). For
\(N=\max\{i,j\}\), nestedness gives \(q_1,q_2\in S_N\), and \(b\) is strongly
convex on \(S_N\).

For completeness, strong convexity yields the quantitative separation
directly. For \(0<t<1\),
\[
b(q_1+t(q_2-q_1))
\leq (1-t)b(q_1)+t b(q_2)
-\frac{\mu_N}{2}t(1-t)\lVert q_2-q_1\rVert^2.
\]
Subtract \(b(q_1)\), divide by \(t\), and let \(t\downarrow0\). Since
\(q_1\in U\), Gateaux differentiability gives
\[
B(q_2,q_1)
\geq \frac{\mu_N}{2}\lVert q_2-q_1\rVert^2.
\]
The cluster-point calculation in the proposition gives
\(B(q_2,q_1)=0\), hence \(q_1=q_2\). Therefore the whole sequence converges
weakly, and its limit belongs to \(\MIN(F)\).
\end{proof}

\begin{remark}
The proof shows a sharper replacement for global strict convexity: it is
enough that the Bregman divergence separate pairs of weak cluster points.
Strong convexity on a nested cover of \(C\) is more than sufficient.
\end{remark}

\section*{Broader interpretation: failure without local separation}

The wording of Remark 6.4 mentions choosing an arbitrary point from a possibly
set-valued proximal map. Under Assumption 5.2 the map is actually singleton,
as observed above. If the intended question also removes Assumption 5.2, the
answer changes.

\begin{theorem}
If local strong convexity (or another separation condition on cluster points)
is deleted together with global strict convexity, arbitrary proximal selections
need not converge weakly, even in dimension two with bounded \(U\) and a zero
objective.
\end{theorem}

\begin{proof}
Let \(X=\R^2\) with its Euclidean norm and define
\[
b(s,t)=
\begin{cases}
\tfrac12s^2,&(s,t)\in[-2,2]^2,\\
+\infty,&\text{otherwise}.
\end{cases}
\]
Then \(b\) is proper, lower semicontinuous, and convex, with bounded zone
\(U=(-2,2)^2\), and \(b'(s,t)=(s,0)\) on \(U\). Its derivative is
weak-to-weak-star sequentially continuous, and
\[
B((a,c),(s,t))=\tfrac12(a-s)^2.
\]

Take \(C=S_k=[-1,1]^2\) for every \(k\), and \(f=g=0\). Then
\(\MIN(F)=C\subset U\), the gradient of \(f\) is Lipschitz, every proximal
subproblem has a nonempty solution set contained in \(U\), and the boundedness
alternative for the weak-convergence hypotheses is automatic because \(U\)
is bounded. For arbitrary positive \(L,\mu\), the proximal subproblem at
\((s,t)\in C\) minimizes
\[
\frac{L}{2\mu}(a-s)^2
\quad\text{over }(a,c)\in C.
\]
Consequently its solution set is
\[
\{(s,c):-1\leq c\leq1\}.
\]
Starting with \(x_1=(0,1)\), choose successively
\(x_k=(0,(-1)^{k+1})\). Every \(x_k\) is an exact allowed proximal choice from
the preceding point, but the sequence alternates between two points. In finite
dimension weak and norm convergence coincide, so it has no weak limit.

This example intentionally violates Assumption 5.2: \(b\) is flat in the
second coordinate and is not strongly convex on \(C\). It therefore does not
contradict the affirmative theorem; it identifies precisely why the retained
local hypothesis matters.
\end{proof}

\section*{Verification notes and limitations}

\begin{itemize}[leftmargin=*]
\item The proof is analytic; no numerical experiment is used as evidence.
\item The source proof of Lemma 8.1 was checked line by line. Its global
strict-convexity invocation occurs only after the limiting-difference
calculation; the local estimate in the main theorem replaces exactly that
step.
\item The evidence crops are reproducible from the packet directory with
\begin{center}
\small\texttt{conda run -n sandbox python}\\
\texttt{code/make\_open\_problem\_crops.py}
\end{center}
\item Human review should focus on the interpretation of Remark 6.4 and on
the fact that Assumption 5.2 is formally retained. Both interpretations are
mathematically covered here.
\item The novelty search was bounded and cannot exclude an unindexed answer.
\end{itemize}

\begin{thebibliography}{9}
\bibitem{RRD}
D. Reem, S. Reich, and A. De Pierro,
\emph{A Telescoping Bregmanian Proximal Gradient Method Without the Global
Lipschitz Continuity Assumption}, Journal of Optimization Theory and
Applications 182 (2019), 851--884; arXiv:1804.10273v4.
\end{thebibliography}

\end{document}
